\section{The shared Boolean and bounded-search kernel}
\label{sec:boolean}

\subsection{Boolean predicates}

A bounded numerical predicate is represented by a primitive-recursive bit
\(b(\bar x)\), with a normalization theorem
\[
  b(\bar x)=0\;\vee\;b(\bar x)=1.
\]
We write \([P](\bar x)\) for the object-language equation
\(b_P(\bar x)=1\).  Boolean conjunction, disjunction, implication, and
negation are fixed truth-table functions.  Equality and order are represented
by fixed Boolean tests.

For a bit predicate \(b(i,\bar x)\), define finite folds
\begin{align*}
\All_b(0,\bar x)&=1,
&\All_b(n+1,\bar x)&=\All_b(n,\bar x)\wedge b(n,\bar x),\\
\Any_b(0,\bar x)&=0,
&\Any_b(n+1,\bar x)&=\Any_b(n,\bar x)\vee b(n,\bar x).
\end{align*}
A least search \(\Search_b(n,\bar x)\) returns the least \(i<n\) with
\(b(i,\bar x)=1\), or the sentinel \(n\) if none exists.

\begin{theorem}[Fold and search correctness]
\label{thm:fold-correct}
For every fixed primitive-recursive bit predicate \(b\), iPRA proves:
\begin{align*}
\All_b(n,\bar x)=1
  &\leftrightarrow \forall i<n\; b(i,\bar x)=1,\\
\Any_b(n,\bar x)=1
  &\leftrightarrow \exists i<n\; b(i,\bar x)=1.
\end{align*}
If \(\Any_b(n,\bar x)=1\), then
\(\Search_b(n,\bar x)<n\), the returned index satisfies \(b=1\), and every
smaller index fails.
\end{theorem}

\begin{proof}
The proof is organized around Boolean equations, not logical quantifiers inside
the induction predicate.

For the universal fold, fix an index \(i\) and induct on \(n\) using the
quantifier-free formula
\[
  [i<n]\wedge [\All_b(n,\bar x)=1]
  \;\longrightarrow\;[b(i,\bar x)=1].
\]
At zero the antecedent is false.  At a successor, the decidable numerical split
\(i<n\) or \(i=n\) reduces the claim respectively to the induction hypothesis
or to the newest factor.  Universal introduction over \(i\) occurs after the
induction.  The reverse implication is proved by induction on the prefix
length, keeping the quantified hypothesis in the context and instantiating it
at the current index.

The existential fold is dual.  The forward proof inspects the last bit; the
reverse proof propagates a supplied witness into all longer prefixes.  The
search routine carries a found bit and a candidate.  Its quantifier-free
invariant states: either no earlier bit is one, or the stored candidate is in
range, satisfies the predicate, and no smaller candidate does.  At termination,
the existential fold excludes the no-witness branch.
\end{proof}

\subsection{Structural Booleanization}

\begin{theorem}[Bounded Booleanization]
\label{thm:booleanization}
Let \(\theta(\bar x)\) be a fixed formula built from equations between
primitive-recursive terms, propositional connectives, and bounded quantifiers
\(\forall i<B(\bar x)\) and \(\exists i<B(\bar x)\).  There is an effectively
constructed \(\PR\) bit symbol \(b_\theta\) such that iPRA proves
\[
 b_\theta(\bar x)=1\quad\leftrightarrow\quad\theta(\bar x).
\]
\end{theorem}

\begin{proof}
Proceed by structural recursion.  Atomic equations use the equality bit.
Propositional connectives use Boolean truth functions; their correctness follows
by the two numerical cases for each input bit and is intuitionistically valid
because the bits are decidable.  A bounded universal becomes an \(\All\)-fold,
and a bounded existential becomes an \(\Any\)-fold.  Apply
\cref{thm:fold-correct}.  Composition and primitive recursion preserve
primitive recursiveness, and the proof derivation is assembled from the
component equivalences by substitution and natural-deduction rules.
\end{proof}

\subsection{Finite minimization and canonicalization}

Let \(A(i,p)\) be a Boolean admissibility predicate and let
\(c(i,p)\) be a coded rational or integer cost with a decidable comparison.
A primitive recursion scans \(i<C(p)\), retains a found bit, and stores the
least-cost candidate, breaking ties by the least index.

\begin{lemma}[Finite minimization]
\label{lem:minimize}
If some candidate below \(C(p)\) is admissible, iPRA proves that the scan
returns an admissible index below \(C(p)\) whose cost is no larger than that of
any admissible candidate; among equal minima it returns the least index.
\end{lemma}

\begin{proof}
At prefix length \(j\), the invariant says either no admissible candidate has
appeared, or the stored index is admissible, lies below \(j\), is cost-minimal
among earlier candidates, and satisfies the tie-breaking clause.  All fields
are Boolean equations or bounded comparisons.  The successor step performs one
admissibility test and one comparison.  At the final prefix, the nonemptiness
premise rules out the first branch.
\end{proof}

Finite minimization is the canonicalization mechanism used throughout the
report.  It does not assert that the chosen code is mathematically canonical;
it merely supplies a deterministic representative from a proved nonempty
finite relation.  Extensionality is established separately.

\subsection{A standard compilation rule}

The recommended rule for every bounded assertion is:
\begin{enumerate}[label=(\arabic*)]
\item define a fixed \(\PR\) checker;
\item prove its local recursion equations;
\item prove a quantifier-free Boolean invariant;
\item derive the ordinary logical statement after induction; and
\item export the resulting derivation, not only a semantic truth theorem.
\end{enumerate}
This rule prevents bounded quantifiers from leaking into every later induction
and makes certificate verification, path checking, finite nets, and search
share one small kernel.
